\input{preamble}
\input{volumes/references/volume-2-signed-kernels}
\IfFileExists{volume-1-polylogarithm-identities.aux}{\externaldocument[][nocite]{volume-1-polylogarithm-identities}[volume-1-polylogarithm-identities.pdf]}{}
\IfFileExists{volume-3-hurwitz-stieltjes.aux}{\externaldocument[][nocite]{volume-3-hurwitz-stieltjes}[volume-3-hurwitz-stieltjes.pdf]}{}
\IfFileExists{volume-4-cm-and-herglotz.aux}{\externaldocument[][nocite]{volume-4-cm-and-herglotz}[volume-4-cm-and-herglotz.pdf]}{}
\title{Polylogarithms and their Arithmetic Bridges\\[8pt]
\large Volume 2: Signed Kernels, Zero Geometry and Experimental Discovery}
\author{ProveIt Contributors}
\date{Collective manuscript, October 2026}
\hypersetup{pdftitle={Polylogarithms and their Arithmetic Bridges, Volume 2: Signed Kernels, Zero Geometry and Experimental Discovery},pdfauthor={ProveIt Contributors}}
\begin{document}
\pagestyle{plain}
\frontmatter
\maketitle
\chapter*{About this volume}
\addcontentsline{toc}{chapter}{About this volume}
Signed moment kernels, exact angular zero counts, all-depth compensation and certified computation develop the analytic geometry. The experimental chapter explains how computations generate proof obligations and how exact certificates settle them.

The four volumes form one collective manuscript. They share definitions,
notation, source evidence and a bibliography. Chapter, theorem and equation
numbers are retained from the unified manuscript so references in research
reports remain stable. Cross-volume references link to the companion PDF;
keep the four PDF files together when downloading them.

\begin{center}
\begin{tabular}{@{}p{0.8in}p{5.5in}@{}}
\href{volume-1-polylogarithm-identities.pdf}{Volume 1} & Polylogarithm Identities and Arithmetic Reductions \\
\href{volume-2-signed-kernels.pdf}{Volume 2} & Signed Kernels, Zero Geometry and Experimental Discovery \\
\href{volume-3-hurwitz-stieltjes.pdf}{Volume 3} & Hurwitz Jets, Stieltjes Calculus and Gamma Correlations \\
\href{volume-4-cm-and-herglotz.pdf}{Volume 4} & Gamma Grids, CM Periods and Herglotz Arithmetic \\
\end{tabular}
\end{center}

A rigorous analytic proof, a finite exact certificate and an independently
repeated numerical diagnostic carry different evidence. Computation drives
conjecture formation and exposes missing structure; successful numerical
recognition does not itself prove an identity or arithmetic independence.
The source inventory and editorial ledger record the accepted corrections
and the scope of material still undergoing analytic reconciliation.

\tableofcontents
\mainmatter
\setcounter{chapter}{4}
\input{chapters/05-certified-computation}
\setcounter{chapter}{11}
\input{chapters/10-discovery}
\appendix
\input{chapters/11-literature}
\backmatter
\input{references}
\end{document}
